\section{Preliminaries}
%%%%%%%%%%%%%%%%%%%%%%%%%%%%%%%%%%%%%%%%%%%%%%

For a Hadamard manifold $M$ and points $\xi\in M_\infty$ and $x\in M$,
the Busemann function $b$ associated to $\xi$ with $b(x)=0$ is given by
\begin{align*}
	b(y) = \lim_{n\to\infty} (d(x_n,y)-d(x_n,x)),
\end{align*}
where $(x_n)$ is any sequence in $M$ converging to $\xi$.
Recall that the limit exists and that $b$ is a $C^2$ distance function on $M$
with horoballs and horospheres centered at $\xi$ as sublevel and level sets;
cf.~\cite{Ba95} for this and other facts about Hadamard manifolds.

\begin{obs}\label{chee}
If a Hadamard manifold $M$ is asymptotically harmonic about a point $\xi\in M_\infty$
or if $N=\Gamma\backslash M$ is of type {\rm \ref{pa}},
then the Cheeger constant of $M$ respectively $N$ is at least $h$,
where~$h$ is the mean curvature of the horospheres with center $\xi$.
In particular, the spectrum of~$M$ respectively $N$ is contained in $\bigl[h^2/4,\infty\bigr)$.
\end{obs}

Throughout the article,
we use the absolute value sign to denote the appropriate volumes of sets under discussion.

\begin{proof}[Proof of Observation~\ref{chee}]
Let $D$ be a compact domain in $M$ respectively $N$ with smooth boundary.
Let $b$ be a Busemann function associated to $\xi$ and $X=\nabla b$.
Then $|X|=1$ and $\diver X=h$ and hence
\begin{align*}
	h|D| = \int_D\diver X = \int_{\partial D}\langle X,\nu\rangle \le |\partial D|,
\end{align*}
where $\nu$ is the outer normal field of $D$ along $\partial D$.
The last assertion is just the Cheeger in\-equality.
\end{proof}

\begin{obs}\label{obs}
If a Hadamard manifold $M$ is asymptotically harmonic about a point $\xi\in M_\infty$
and $b$ is a Busemann function of $M$ centered at $\xi$, then
\begin{align*}
	0 \le \nabla^2b(v,v) \le h |v|^2.
\end{align*}
\end{obs}

\begin{proof}
The claim follows immediately from $h=\tr\nabla^2b$ and $\nabla^2b\ge0$.
\end{proof}
